Els ratpenats emeten uns xiscles en forma d'ultrasons i utilitzen els ecos d'aquests ultrasons per a orientar-se i per a detectar obstacles i preses. Una espècie de ratpenats emet ultrasons amb una freqüència de 83,0~kHz quan caça mosquits.

\figura[0.4]{fig1.png}

\begin{apartat}{1}
Calculeu la longitud d'ona i el període dels ultrasons emesos per aquests ratpenats. Considereu un mosquit situat a 1,5000~m de l'orella dreta i a 1,5030~m de l'orella esquerra del ratpenat. Calculeu la diferència de fase en l'eco percebut per cada orella, provinent del mosquit.
\end{apartat}

\begin{apartat}{1}
Quan el mosquit està més a prop, el ratpenat també podria utilitzar la diferència d'intensitats dels ecos. Calculeu el quocient d'intensitats sonores $\dfrac{I_\mathrm{dreta}}{I_\mathrm{esquerra}}$ quan el mosquit està a 33~cm de l'orella dreta i a 34~cm de l'orella esquerra i expresseu en decibels la diferència de nivells d'intensitat sonora. Considereu que l'eco es propaga uniformement des del mosquit en totes les direccions de l'espai.
\end{apartat}

\blocetiquetat{Dada:}{Velocitat dels ultrasons en l'aire $= 340\ \mathrm{m\,s^{-1}}$.}
